\scanpage{15}
\mspage{12}

Orientation du segment : choix d'une des extrémités, ou encore, d'un des deux ordres qui \unsure{définit} la structure de segment.

\begin{center}\rule{3cm}{0.4pt}\end{center}

[Application aux modèles réguliers (i.e. tels que $L=L_r$)]. Considérons la relation sur $L$, en $x,y$
\[
\exists I\in S\text{ tel que }x,y\in\tilde I
\]
[en vertu du cor. 2 à F4), \inserted{si $L=L_r$, est} \unsure{équivalent} à
\[
\exists I\in S\text{ tel que }x,y\in\tilde I^\circ.
\]
]

Cette relation est réflexive et symétrique. Considérons la relation d'équivalence engendrée. L'ens des classes d'équivalence est noté $\pi_0(\mathcal X)$ (composantes connexes). On a une notion évidente de sommes directes des modèles de $1$ formes, et on voit que \unsure{tte} $1$-forme s'identifie

\mspage{14}

\rotatebox{90}{\reflectbox{\unsure{EXTRA STRONG}}}
\rotatebox{90}{\reflectbox{\unsure{EXTRA STRONG}}}
